% !TeX program = lualatex
\documentclass[11pt]{ltjsarticle}
\usepackage{amsmath,amssymb,amsthm,mathtools}
\usepackage{lmodern}
\usepackage[unicode]{hyperref}

\theoremstyle{plain}
\newtheorem{theorem}{定理}[section]
\newtheorem{proposition}[theorem]{命題}
\newtheorem{lemma}[theorem]{補題}
\newtheorem{corollary}[theorem]{系}

\theoremstyle{definition}
\newtheorem{definition}[theorem]{定義}
\newtheorem{assumption}[theorem]{仮定}
\newtheorem{remark}[theorem]{注意}

\newcommand{\PrT}{\mathrm{Pr}_T}
\newcommand{\SelfRef}{\mathrm{SelfRef}}
\newcommand{\QRB}{\mathrm{QRB}}
\newcommand{\Res}{\mathrm{Res}}
\newcommand{\Supp}{\operatorname{supp}}
\newcommand{\Spc}{\operatorname{Spc}}
\newcommand{\dfix}{\bigcirc}

\title{言語の外点：\\
---強制法と四重残余束による到達不能性の幾何学---}
\author{千葉将臣}
\date{2026年7月19日}

\begin{document}

\maketitle

\begin{abstract}
本論文は、形式的体系の内部では対象として定義できないが、内部からの関係によって間接的に指示できるものを「外点」\(\dfix\) と呼び、その数学的記述を試みる。
第一に、完備ブール代数値モデルと強制関係を用い、基底モデルに属さないジェネリック・フィルターを、基底モデル内で定義可能な関係によって指示する。
第二に、四種の操作に伴う情報損失を安定圏における余ファイバーとして定式化し、それらの台の共通部分を四重残余束（Quadruple Residual Bundle; QRB）の特異領域と定義する。
強制法に関する標準定理とQRBに固有の仮定を明確に分離することで、外点とジェネリック・フィルターの同一視を無条件の定理とはせず、比較写像が満たすべき条件として提示する。
これにより、到達不能性を単なる未定義性ではなく、複数の比較写像が同時に同値でなくなる幾何学的障害として研究するための形式的枠組みを与える。
\end{abstract}

\section{はじめに}

\subsection{問題設定}

十分に強い形式的体系は、自身の構文を算術化できる一方、自身の真理を内部で完全に定義することはできない\cite{godel1931,tarski1936}。
Kripke の真理理論は部分的真理述語の最小不動点を構成するが、得られた言語の外側から新たな真理概念を問うリベンジ問題は残る\cite{kripke1975}。
したがって、本論文では外部を新しい対象定数として内部言語に追加するのではなく、内部で計算できる関係が外部に対して与える制約によって指示する。

この方針には二つの段階がある。
第一段階は、強制関係 \(p\Vdash\varphi\) による外部の指示である。
第二段階は、複数の論理的・圏論的操作が捨てる情報を残余として記録し、その共通の非自明領域を幾何学化することである。

\subsection{本論文の寄与と射程}

本論文の寄与は次の三点である。
\begin{enumerate}
\item 強制法の定義可能性補題と真理補題を、外部の「定義」ではなく「指示」のモデルとして整理する。
\item 安定テンソル三角圏において四つの残余対象とQRBを定義し、その特異領域を台の交差として与える。
\item ジェネリック・フィルターとQRB特異領域の対応を、無条件の同一視ではなく、明示的な比較写像を要する研究仮説として定式化する。
\end{enumerate}

特に、強制法が任意の集合論的命題を保存するとは主張しない。
強制拡大は連続体仮説のような命題を変え得るが、基底モデルが推移的である場合には順序数を保存し、同じ自然数構造上の一階算術文の真理を保存する。

\section{形式的準備}
\label{sec:prelim}

\subsection{体系と自己言及}

\begin{assumption}\label{ass:T}
\(T\) は再帰的に公理化された無矛盾な一階理論であり、少なくとも Robinson 算術 \(Q\) を含み、自身の構文と証明関係を算術化できるとする。
\end{assumption}

論理式 \(\varphi\) の Gödel 数を \(\ulcorner\varphi\urcorner\)、\(T\) の標準的な証明可能性述語を \(\PrT(x)\) と書く。

\begin{lemma}[対角線補題]\label{lem:diagonal}
任意の一変数論理式 \(\psi(x)\) に対して、文 \(G\) が存在し、
\[
T\vdash G\leftrightarrow\psi(\ulcorner G\urcorner)
\]
が成り立つ。
\end{lemma}

補題\ref{lem:diagonal}で \(\psi(x)\equiv\neg\PrT(x)\) とすれば、適切な追加条件の下で Gödel 文が得られる\cite{kleene1952}。
記法として
\[
\SelfRef_T(\varphi)=\neg\PrT(\ulcorner\varphi\urcorner)
\]
と置く。ただし、これは対象言語の全真理を与える作用素ではない。

\subsection{トポスと二重否定層化}

\begin{definition}[エレメンタリー・トポス]
圏 \(\mathcal{E}\) がエレメンタリー・トポスであるとは、有限極限を持ち、デカルト閉であり、部分対象分類子 \(\Omega\) を持つことをいう。
\end{definition}

一般のトポスの内部論理は直観主義的である。
Lawvere--Tierney 位相 \(j=\neg\neg\) に伴う層化関手 \(a_{\neg\neg}\) は、同型を除いて冪等である\cite{johnstone2002}。
したがって、二重否定層化から生じる残余は「層化が非冪等であること」ではなく、単位写像
\(X\to a_{\neg\neg}X\) が一般には同型でないことによって測るべきである。

\section{強制法による外部の指示}
\label{sec:forcing}

\subsection{ブール値と強制関係}

\begin{assumption}\label{ass:forcing}
\(M\) を集合論の十分な有限部分を満たす可算推移モデル、\(\mathbb{P}\in M\) を分離的な強制半順序、\(\mathbb{B}\) を \(M\) におけるその正則開完備化とする。
\end{assumption}

\(M^{\mathbb{B}}\) を \(\mathbb{B}\) 値宇宙とし、論理式 \(\varphi\) のブール値を
\(\lVert\varphi\rVert_{\mathbb{B}}\in\mathbb{B}\) と書く。

\begin{definition}[強制関係]
\(p\in\mathbb{P}\) と強制言語の論理式 \(\varphi\) に対して
\[
p\Vdash\varphi
\quad\Longleftrightarrow\quad
p\leq\lVert\varphi\rVert_{\mathbb{B}}
\]
と定める。
\end{definition}

\begin{theorem}[定義可能性補題と真理補題]\label{thm:forcing}
仮定\ref{ass:forcing}の下で、各論理式 \(\varphi\) に対する関係
\(p\Vdash\varphi\) は \(M\) 内で再帰的に定義できる。
さらに、\(G\subseteq\mathbb{P}\) が \(M\)-ジェネリックならば
\[
M[G]\models\varphi(\tau_1^G,\ldots,\tau_n^G)
\quad\Longleftrightarrow\quad
(\exists p\in G)\;p\Vdash\varphi(\tau_1,\ldots,\tau_n)
\]
が成り立つ。
\end{theorem}

定理\ref{thm:forcing}は強制法の標準的な基本定理である\cite{cohen1963,bell1985,jech2003}。
右辺は \(M\) 内で扱えるが、ジェネリック・フィルター \(G\) 自身は一般に \(M\) の元ではない。
この非対称性が、本論文における「外部を定義せずに指示する」という考えの厳密な原型である。

\begin{definition}[強制法における外点]\label{def:forcing-outer-point}
基底モデル \(M\) の外部にある \(M\)-ジェネリック・フィルター \(G\) を、強制法によって指示される外点 \(\dfix\) の意味論的代表と呼び、
\[
\dfix_{\mathrm{Forc}}\coloneqq G
\]
と表す。
この等式は \(M\) の内部で対象 \(G\) を定義するものではなく、メタ理論において外点 \(\dfix\) の役割を \(G\) に担わせるという意味論的な同定である。
体系内部から利用できるのは \(G\) 自体ではなく、定理\ref{thm:forcing}の強制関係 \(p\Vdash\varphi\) を通じて記述される \(\dfix_{\mathrm{Forc}}\) の条件付きの振る舞いである。
\end{definition}

\begin{proposition}[相対的無矛盾性と算術的絶対性]\label{prop:preservation}
\(M\models\mathrm{ZFC}\) かつ \(G\) が \(M\)-ジェネリックならば、\(M[G]\models\mathrm{ZFC}\) であり、\(M\) と \(M[G]\) は同じ順序数を持つ。
さらに両者が標準的な \(\omega\) を共有する限り、自然数構造だけを量化領域とする一階算術文の真理は一致する。
\end{proposition}

\begin{remark}
命題\ref{prop:preservation}は、強制拡大が集合論の全命題について保存拡大であることを意味しない。
新しい実数、すなわち \(\omega\) の新しい部分集合が追加される場合があり、集合論的命題の真理値は変化し得る。
\end{remark}

\section{四重残余束の構成}
\label{sec:qrb}

\subsection{残余対象}

QRBを文字通りの位相的ベクトル束と仮定するのではなく、まず余ファイバーを持つ圏で残余を定義する。

\begin{assumption}\label{ass:category}
\(\mathcal{C}\) を冪等完備な安定テンソル三角圏とし、\(X\in\mathcal{C}^c\) を体系 \(T\) の意味論的状態を表すコンパクト対象とする。
各 \(i\in\{S,N,SN,NS\}\) に対し、完全自己関手 \(M_i\colon\mathcal{C}\to\mathcal{C}\) と自然変換
\(\eta_i\colon\mathrm{Id}_{\mathcal{C}}\Rightarrow M_i\) が与えられ、\(M_i\) はコンパクト対象を保つとする。
\end{assumption}

\begin{definition}[四つの残余]
各操作の残余対象を完全三角
\[
X\xrightarrow{\eta_i(X)}M_iX\longrightarrow\Res_i(X)\longrightarrow\Sigma X
\]
によって定まる余ファイバー
\[
\Res_i(X)=\operatorname{cofib}\!\bigl(\eta_i(X)\bigr)
\]
と定義する。
\end{definition}

想定する四操作は次の通りである。
\begin{enumerate}
\item \(M_S\)：同一性を外延的等号へ截断する比較操作。単なる恒等関手を選べば残余はゼロになるため、非自明な比較写像を明示する必要がある。
\item \(M_N\)：二重否定層化。残余 \(\Res_N(X)\) は \(X\to a_{\neg\neg}X\) が同型でない度合いを表す。
\item \(M_{SN}\)：随伴 \(F\dashv U\) から生じるモナド \(UF\)。残余は単位 \(X\to UFX\) の非可逆性を表す。
\item \(M_{NS}\)：構文の対角化を意味論側に実現する自己参照比較操作。この実現の存在は追加構造であり、対角線補題だけから自動的には得られない。
\end{enumerate}

\begin{definition}[四重残余束]
対象 \(X\) における四重残余束を
\[
\QRB(X)=
\Res_S(X)\oplus\Res_N(X)\oplus
\Res_{SN}(X)\oplus\Res_{NS}(X)
\]
と定義する。
\end{definition}

\subsection{台と外点領域}

\(\mathcal{C}^c\) を \(\mathcal{C}\) のコンパクト対象の部分圏とし、Balmer スペクトル
\(\Spc(\mathcal{C}^c)\) 上の台を \(\Supp(-)\) と書く\cite{balmer2005}。

\begin{definition}[外点領域]
四つの残余が同時に非自明となる領域を
\[
\mathcal{O}_X=
\bigcap_{i\in\{S,N,SN,NS\}}\Supp\bigl(\Res_i(X)\bigr)
\subseteq\Spc(\mathcal{C}^c)
\]
と定義する。
\(\mathcal{O}_X\neq\varnothing\) のとき、その点をQRBの外点候補と呼ぶ。
\end{definition}

ここで \(\mathcal O_X\) は外点 \(\dfix\) そのものではなく、四つの残余が同時に検出される外点候補の領域である。
したがって、強制法における代表 \(\dfix_{\mathrm{Forc}}=G\) と対応させるべき対象は領域全体 \(\mathcal O_X\) ではなく、その内部の特定の点である。
また、\(\dfix\)は \(\mathcal O_X\) の元（点）であり、\(\mathcal O_X\) はその候補領域である。

\begin{theorem}[交差点の判定]\label{thm:intersection}
仮定\ref{ass:category}の下で、素点 \(\mathfrak{p}\in\Spc(\mathcal{C}^c)\) が \(\mathcal{O}_X\) に属することと、すべての \(i\) について局所化された比較写像
\[
\eta_i(X)_{\mathfrak p}\colon
X_{\mathfrak p}\longrightarrow(M_iX)_{\mathfrak p}
\]
が同型でないことは同値である。
\end{theorem}

\begin{proof}
安定圏では射が同型であることとその余ファイバーがゼロ対象であることが同値である。
また、\(\mathfrak p\notin\Supp(Y)\) であることは局所化 \(Y_{\mathfrak p}\) がゼロであることと同値である。
これを四つの余ファイバーへ適用すればよい。
\end{proof}

定理\ref{thm:intersection}により、外点領域は比喩ではなく、四つの比較写像が同時に可逆性を失う素点の集合として特徴づけられる。
ただし、その非空性は一般には保証されず、具体的な圏と四操作ごとに証明されなければならない。

\section{強制法とQRBを結ぶ比較原理}
\label{sec:comparison}

ジェネリック・フィルター \(G\) と \(\mathcal{O}_X\) は、ここまでの定義だけでは異なる種類の対象である。
前者は強制半順序上のフィルターであり、後者はテンソル三角スペクトルの部分集合である。
両者を同一視するには比較写像が必要となる。

\begin{assumption}[指示写像]\label{ass:designation}
順序を保つ写像
\[
\Theta\colon\mathbb{B}\longrightarrow
\operatorname{Open}\!\bigl(\Spc(\mathcal{C}^c)\bigr)
\]
が存在し、有限交叉と任意結合を保存するとする。
さらに、点 \(\xi_G\in\mathcal O_X\) が存在し、各 \(p\in\mathbb P\) に対して
\(p\in G\) と \(\xi_G\in\Theta(p)\) が同値であるとする。
\end{assumption}

\begin{definition}[QRBにおける外点]\label{def:qrb-outer-point}
仮定\ref{ass:designation}を満たす点 \(\xi_G\in\mathcal O_X\) を、QRBによって表現された外点と呼び、
\[
\dfix_{\mathrm{QRB}}\coloneqq\xi_G
\]
と定める。
このとき外点の二つの表現の関係を
\[
\dfix
\quad\longleftrightarrow\quad
\dfix_{\mathrm{Forc}}=G
\quad\longleftrightarrow\quad
\dfix_{\mathrm{QRB}}=\xi_G\in\mathcal O_X
\]
と記す。
右側の対応は仮定\ref{ass:designation}に依存し、\(G\) と \(\xi_G\) が同一の集合論的対象であることではなく、条件 \(p\) の所属と開集合 \(\Theta(p)\) への所属が一致するという指示構造の同値を意味する。
\end{definition}

\begin{proposition}[条件付き指示原理]\label{prop:designation}
仮定\ref{ass:designation}の下で、\(p\Vdash\varphi\) ならば
\[
\mathcal{O}_X\cap\Theta(p)
\subseteq
\mathcal{O}_X\cap\Theta\!\left(\lVert\varphi\rVert_{\mathbb B}\right)
\]
が成り立つ。
したがって、QRB外点領域上でも強制関係と整合する包含関係によって \(\varphi\) を指示できる。
\end{proposition}

\begin{proof}
\(p\Vdash\varphi\) の定義から \(p\leq\lVert\varphi\rVert_{\mathbb B}\) である。
\(\Theta\) の単調性により
\(\Theta(p)\subseteq\Theta(\lVert\varphi\rVert_{\mathbb B})\) を得る。
両辺と \(\mathcal{O}_X\) の交わりを取れば結論が従う。
\end{proof}

命題\ref{prop:designation}が与えるのは、ジェネリック・フィルターとQRBの完全な同型ではなく、両者の指示構造が整合するための十分条件である。
すなわち、体系内部から強制関係によって指示される外点 \(\dfix_{\mathrm{Forc}}\) の振る舞いが、QRB側では点 \(\dfix_{\mathrm{QRB}}\) の開近傍への所属関係として再表現される。
具体例を構成するには、\(\mathbb B\)、\(\mathcal C\)、四関手、および \(\Theta\) を同時に指定し、仮定\ref{ass:designation}を検証する必要がある。

\section{結論}
\label{sec:conclusion}

本論文は、言語の外点を二つの水準に分けて定式化した。
強制法の水準では、基底モデルに属さないジェネリック・フィルターを \(\dfix_{\mathrm{Forc}}=G\) として、内部で定義可能な強制関係によって間接的に指示できる。
圏論的水準では、四つの比較操作の余ファイバーをQRBとして束ね、その台の共通部分を外点候補領域 \(\mathcal O_X\) とし、指示写像が存在する場合には \(\dfix_{\mathrm{QRB}}=\xi_G\in\mathcal O_X\) と定めた。

この分離によって、標準的な強制定理、圏論から従う交差点判定、そして今後構成すべき指示写像を区別できる。
今後の主要課題は、具体的な論理トポスまたは安定圏において \(\mathcal O_X\neq\varnothing\) を示し、強制代数から台の開集合への写像 \(\Theta\) を構成することである。
その達成によって初めて、ジェネリックな外部と四重残余の特異領域との対応が、比喩ではなく具体的な数学的定理となる。

\section*{謝辞}
本稿の構想に関する議論に感謝する。

\begin{thebibliography}{99}

\bibitem{balmer2005}
Paul Balmer,
``The spectrum of prime ideals in tensor triangulated categories,''
\emph{Journal für die reine und angewandte Mathematik}, 588, 149--168, 2005.

\bibitem{bell1985}
John L. Bell,
\emph{Boolean-Valued Models and Independence Proofs in Set Theory},
Oxford University Press, 1985.

\bibitem{cohen1963}
Paul J. Cohen,
``The independence of the continuum hypothesis,''
\emph{Proceedings of the National Academy of Sciences}, 50, 1143--1148, 1963.

\bibitem{godel1931}
Kurt Gödel,
``Über formal unentscheidbare Sätze der Principia Mathematica und verwandter Systeme I,''
\emph{Monatshefte für Mathematik und Physik}, 38, 173--198, 1931.

\bibitem{jech2003}
Thomas Jech,
\emph{Set Theory: The Third Millennium Edition, Revised and Expanded},
Springer, 2003.

\bibitem{johnstone2002}
Peter T. Johnstone,
\emph{Sketches of an Elephant: A Topos Theory Compendium},
Oxford University Press, 2002.

\bibitem{kleene1952}
Stephen C. Kleene,
\emph{Introduction to Metamathematics},
North-Holland, 1952.

\bibitem{kripke1975}
Saul A. Kripke,
``Outline of a theory of truth,''
\emph{The Journal of Philosophy}, 72(19), 690--716, 1975.

\bibitem{tarski1936}
Alfred Tarski,
``Der Wahrheitsbegriff in den formalisierten Sprachen,''
\emph{Studia Philosophica}, 1, 261--405, 1936.

\end{thebibliography}

\end{document}
